\scanpage{36}

aux homomorphismes

\[
T\longrightarrow A^{\handwritten{\natural}}\longrightarrow A\ ,
\]

nous montre alors aussitôt que $\underline{F}(A^{\handwritten{\natural}})$ et $\underline{F}(T)$ s'identifient respectivement à

\[
\tag*{(13.4.10)}
\underline{F}(A^{\handwritten{\natural}})\simeq\underline{F}^{f}\overset{\mathrm{dfn}}{=}\underline{F}\cap\mathcal{E}^{f}\ ,\quad \underline{F}(T)=\underline{F}^{t}=\underline{F}\cap\mathcal{E}^{t}=0\ ,
\]

le $=0$ provenant du fait immédiat que pour un tore $T$, on a $\underline{F}(T)=0$, la suite exacte du type (13.4.3) devenant ici $0\to0\to\mathcal{E}(T)\to\underline{\mathrm{Lie}}(T)\to0$.

On voit donc en même temps que l'homomorphisme composé

\[
\mathcal{E}^{t}\longrightarrow\mathcal{E}\longrightarrow\underline{E}
\]

induit un isomorphisme

\handwritten{\mspage{495}}

\mspage{495}

% archive sheet number: 35
